\documentclass[a4paper]{article}

\usepackage[a4paper,margin=25mm]{geometry}
\usepackage[danish]{babel}
\usepackage{titlecaps}

\Addlcwords{a an the and but for or nor as at by for in of on per to via vs off up out}

\usepackage{parskip} 
\usepackage{microtype}
\usepackage{setspace}
\setstretch{1.2}

\usepackage{mathtools, amssymb}
\usepackage{bm}
\renewcommand{\Vec}[1]{\boldsymbol{#1}}
\let\oldhat\hat
\renewcommand{\hat}[1]{\oldhat{\mathbf{#1}}}

\DeclareMathOperator{\grad}{grad}

\usepackage{fourier}

\usepackage[dvipsnames]{xcolor}
\definecolor{AUblue}{RGB}{0,61,115}
\definecolor{AUblueDark}{RGB}{0,37,70}


\usepackage{graphicx}
\graphicspath{{./figures/}}
\usepackage{tikz}
\usetikzlibrary{calc, arrows.meta}

\usepackage{enumitem}
\setlist{topsep=4pt,itemsep=2pt,parsep=0pt}

\usepackage{siunitx}
\sisetup{
  exponent-product = \cdot,
  output-decimal-marker = {,},
  per-mode = fraction,
  round-mode = figures,
  round-precision = 3,
  round-pad = false,
}

%Giles Castelles incfig
\usepackage{import}
\usepackage{xifthen}
\usepackage{pdfpages}
\usepackage{transparent}

\newcommand{\incfig}[2][1]{%
  \def\svgwidth{#1\columnwidth}
  \import{./figures/}{#2.pdf_tex}
}

\usepackage{fancyhdr}
\pagestyle{fancy}
\fancyhf{}
\renewcommand{\headrule}{\hbox to\headwidth{\color{AUblueDark}\leaders\hrule height 0.6pt\hfill}}
\setlength{\headheight}{20pt}
\fancyhead[L]{\small\AuthorName}
\fancyhead[C]{\small\ReportTitle}
\fancyhead[R]{\small\TheDate}
\usepackage{lastpage}
\fancyfoot[R]{\small p. \thepage{} of \pageref*{LastPage}}
\fancyfoot[L]{\small\CourseName}

\newcommand{\NoteType}{Forelæsningsnoter}
\newcommand{\Dept}{Institut for Mekanik og Produktion}
\newcommand{\Uni}{Aarhus Universitet}
\newcommand{\ReportTitle}{\titlecap{\NoteType}}
\newcommand{\AuthorName}{Noah Rahbek Bigum Hansen}
\newcommand{\StudentID}{202405538}
\newcommand{\CourseName}{Maskinelementer}
\newcommand{\CourseNumber}{290211U009}
\newcommand{\Course}{\CourseName{} --- \CourseNumber}
\newcommand{\TheDate}{\today}
\newcommand{\Logo}{figures/AU}


\usepackage{hyperref}
\hypersetup{
  colorlinks=true,
  linkcolor=AUblue,
  citecolor=AUblue,
  urlcolor=AUblue,
  pdftitle={\NoteType},
  pdfauthor={\AuthorName},
  pdfsubject={Lecture Notes by \AuthorName{} (\StudentID) for the course \CourseName{} taught at the \Dept{} of \Uni. Last updated on \TheDate{}.},
  pdfcreator={pdfTeX + newtx},
}
\usepackage[nameinlink]{cleveref}

\usepackage{chngcntr}
\counterwithin{figure}{section}
\renewcommand{\thefigure}{\thesection.\arabic{figure}}

\crefname{pluralequation}{equations}{equations}
\Crefname{pluralequation}{Equations}{Equations}

\usepackage{caption}
\captionsetup{font=small}
\usepackage{subcaption} 

\newenvironment{problem}[1]%
{%
  \par\vspace{0pt}%
  {\large\bfseries\sffamily\color{AUblueDark}%
    Problem #1%
  }%
  \par\vspace{.35\baselineskip}%
}%
{%
  \par\vspace{.6\baselineskip}%
  {\color{AUblueDark}\rule{\linewidth}{0.8pt}}%
  \vspace{.6\baselineskip}\par%
}

\newcommand{\AUheading}[1]{%
  \par\vspace{.4\baselineskip}%
  {\sffamily\bfseries\color{AUblueDark}#1\par}%
  \nobreak%
}

\newenvironment{assumptions}[1][Assumptions]{%
  \AUheading{#1}
  \begin{itemize}[leftmargin=*,nosep]
}{\end{itemize}\par\medskip}

\newenvironment{derivation}[1][Derivation]{%
  \AUheading{#1}
  \begingroup\allowdisplaybreaks
  \setlength{\abovedisplayskip}{6pt}%
  \setlength{\belowdisplayskip}{6pt}%
  \setlength{\abovedisplayshortskip}{4pt}%
  \setlength{\belowdisplayshortskip}{4pt}%
}{\par\endgroup\medskip}

\newcommand{\finalresult}[1]{%
  \vspace{.2\baselineskip}%
  \begin{equation}
    {\color{AUblueDark}\boxed{\displaystyle #1}}
    \tag*{\textsc{\textcolor{AUblueDark}{Result}}}
  \end{equation}
  \newpage
}

\usepackage{environ}
\newlength{\pbfigsep}
\setlength{\pbfigsep}{8mm}

\NewEnviron{problemwithimage}[3][0.38\linewidth]{%
  \begin{problem}{#3}%
    \noindent
    \begin{minipage}[t]{\dimexpr\linewidth-#1-\pbfigsep\relax}
      \vspace{-0.2\baselineskip}\BODY
    \end{minipage}\hfill
    \begin{minipage}[t]{#1}
      \vspace{-\baselineskip}\centering
      \includegraphics[width=\linewidth]{#2}%
    \end{minipage}%
  \end{problem}%
}

\usepackage{thmtools}
\usepackage{tcolorbox}
  \tcbuselibrary{skins, breakable}
  \tcbset{
    space to upper=1em,
    space to lower=1em,
  }

\newtcolorbox[auto counter]{definition}[1][]{%
  breakable,
  colframe=ForestGreen,  %frame color
  colback=ForestGreen!5, %background color
  colbacktitle=ForestGreen!25, %background color for title
  coltitle=ForestGreen!70!black,  %title color
  fonttitle=\bfseries\sffamily, %title font
  left=1em,              %space on left side in box,
  enhanced,              %more options
  frame hidden,          %hide frame
  borderline west={2pt}{0pt}{ForestGreen},  %display left line
  title=Definition \thetcbcounter: #1,
}

\newtcolorbox{greenline}{%
  breakable,
  colframe=ForestGreen,  %frame color
  colback=white,          %remove background color
  left=1em,              %space on left side in box
  enhanced,              %more options
  frame hidden,          %hide frame
  borderline west={2pt}{0pt}{ForestGreen},  %display left line
}

\newtcolorbox[auto counter, number within=section]{dis}[1][]{%
  breakable,
  colframe=NavyBlue,  %frame color
  colback=NavyBlue!5, %background color
  colbacktitle=NavyBlue!25,    %background color for title
  coltitle=NavyBlue!70!black,  %title color
  fonttitle=\bfseries\sffamily, %title font
  left=1em,            %space on left side in box,
  enhanced,            %more options
  frame hidden,        %hide frame
  borderline west={2pt}{0pt}{NavyBlue},  %display left line
  title=Discussion \thetcbcounter: #1
}

\newtcolorbox{blueline}{%
  breakable,
  colframe=NavyBlue,     %frame color
  colback=white,         %remove background
  left=1em,              %space on left side in box,
  enhanced,              %more options
  frame hidden,          %hide frame
  borderline west={2pt}{0pt}{NavyBlue},  %display left line
}

\newtcolorbox{exa}[1][]{%
  breakable,
  colframe=RawSienna,  %frame color
  colback=RawSienna!5, %background color
  colbacktitle=RawSienna!25,    %background color for title
  coltitle=RawSienna!70!black,  %title color
  fonttitle=\bfseries\sffamily, %title font
  left=1em,              %space on left side in box,
  enhanced,              %more options
  frame hidden,          %hide frame
  borderline west={2pt}{0pt}{RawSienna},  %display left line
  title=Example: #1,
}

\newtcolorbox[auto counter, number within=section]{sæt}[1][]{%
  breakable,
  colframe=RawSienna,  %frame color
  colback=RawSienna!5, %background color
  colbacktitle=RawSienna!25,    %background color for title
  coltitle=RawSienna!70!black,  %title color
  fonttitle=\bfseries\sffamily, %title font
  left=1em,              %space on left side in box,
  enhanced,              %more options
  frame hidden,          %hide frame
  borderline west={2pt}{0pt}{RawSienna},  %display left line
  title=Sætning \thetcbcounter: #1,
  before lower={\textbf{Bevis:}\par\vspace{0.5em}},
  colbacklower=RawSienna!25,
}

\newtcolorbox{redline}{%
  breakable,
  colframe=RawSienna,  %frame color
  colback=white,       %Remove background color
  left=1em,            %space on left side in box,
  enhanced,            %more options
  frame hidden,        %hide frame
  borderline west={2pt}{0pt}{RawSienna},  %display left line
}

\newtcolorbox{des}[1][]{%
  breakable,
  colframe=NavyBlue,  %frame color
  colback=NavyBlue!5, %background color
  colbacktitle=NavyBlue!25,    %background color for title
  coltitle=NavyBlue!70!black,  %title color
  fonttitle=\bfseries\sffamily, %title font
  left=1em,              %space on left side in box,
  enhanced,              %more options
  frame hidden,          %hide frame
  borderline west={2pt}{0pt}{NavyBlue},  %display left line
  title=Description of #1,
}

\makeatother
\def\@lecture{}%
\newcommand{\lecture}[3]{
  \ifthenelse{\isempty{#3}}{%
    \def\@lecture{Lecture #1}%
  }{%
    \def\@lecture{Lecture #1: #3}%
  }%
  \refstepcounter{section}%
  \phantomsection
  % Typeset like a subsection (no numbering in the heading)
  \subsection*{\makebox[\textwidth][l]{\@lecture \hfill \color{AUblueDark}\normalfont\small{#2}}}
  % Add a section-level TOC entry with the section number
  \addcontentsline{toc}{section}{\protect\numberline{\thesection}\@lecture}%
}

\makeatletter

\newcommand{\exercise}[1]{%
 \def\@exercise{#1}%
 \subsection{Exercise #1}
}

\makeatother
